\subsection{Revisiting Smart Building Objectives Through Affective Interaction}
\label{sec:ch3-objectives}
We now discuss four shifts in smart building research objectives---design, evaluative, methodological, and knowledge---to better support affective, situated, and interpretive experiences in smart environments (\textbf{RQ2.0}).

\subsubsection{\textbf{Design Objectives: Enabling Interpretive Participation}}
Smart building research prioritises \textit{optimal environmental or affective conditions} where the user is the recipient of improved environmental or emotional conditions. Our findings suggest that interaction with the environment is already interpretive, procedural, and affectively driven. Participants tested whether airflow responded to movement, inferred lighting patterns based on system logic, and decided when to endure or relocate. This is \textit{participation}---not in the sense of participatory co-design (e.g. works identified in a recent scoping review ~\cite{rao2025designing}), but as a lived, affective way of involvement in spatial experience. Participation here develops gradually, shaped by repetition, social context, and bodily attunement. Yet this participation is fragile, not guaranteed, and created under the right conditions, faltering when interpretive effort is unmet.

Thus, we propose a shift in design objectives: \textit{from delivering optimal conditions to sustaining interpretive participation}. This reorients the field’s aim towards supporting occupants’ capacity to interpret, adapt to, and influence their environment. The objective is not just to allow interaction, but to maintain bodily, spatial, and social conditions that make affective engagement possible.

\subsubsection{\textbf{Methodological Objectives: Situated Methodology}}
Smart building research infers experience through extractive approaches and decontextualised data points (e.g.~\cite{daissaoui2020iot, dalton2020habitech}) to derive generalisable insights. This reduces experience to measurable proxies, overlooking the contexts in which it unfolds. In contrast, we draw from interpretive HCI, particularly those shaped by AfI~\cite{ahmadpour2025affective}, which foreground situated and affective dimensions of experience. Our walk-throughs~\cite{crivellaro2015contesting, mitchell2020no} elicited interpretation--- participants voiced ambiguity, spatial metaphors, and ``local vernacular'' that standardised metrics miss. Unlike controlled-condition studies~\cite{bower2019impact}, our methodology enabled a socially embedded lens supporting movement, reflection, conversation, temporal patterns and co-regulation.

This points to a needed shift in smart buildings' methodological objectives: \textit{from sensing and extracting behavioural data to situated and interpretive depth}. This demands methodological pluralism: hybrid approaches combining sensing with narrative, embodied, and reflexive techniques~\cite{moran2013using, bower2019impact}. Such depth is not about sensor resolution or data volume, but about understanding how people interpret and respond to environments over time.

\subsubsection{\textbf{Evaluation Objectives: Supporting Affective Fit}}
Smart building research also establishes measurable comfort, affect~\cite{bower2019impact, moran2013using}, privacy~\cite{schnadelbach2019adaptive}, or agency~\cite{mitchell2020no}. Yet our findings show these are not fixed outcomes, but emerge through fleeting, embodied impressions of \textit{``affective fit''}---whether a space feels “good enough” in the moment. We propose a shift from \textit{measuring comfort} to \textit{supporting affective fit}. We define affective fit as the evolving, somatic, and interpretive sense that a space supports one’s activities, emotions, and intentions, however imperfectly. Unlike comfort, it is a continuous negotiation shaped by task, social context, and time. It may be surfaced through multi-modal reflections (e.g. body maps~\cite{turmo2023towards}, narrative probes~\cite{gaver2004cultural}), relational indicators (e.g. space dwell patterns~\cite{verma2017studying}), or technologies that register stability and re-alignment.

Supporting affective fit requires expanding methodologies. Embodied and interpretive methods such as body mapping, which we found effective for surfacing tacit judgments, should be integrated with sensed and behavioural data. Evaluation should enable occupants to interpret their own experience in context and ensure these expressions meaningfully inform building interaction and design.

\subsubsection{\textbf{Knowledge Objectives: Developing Situated Understandings}}

Another objective has been to derive predictive models that hold across contexts and users. This logic underpins optimisation frameworks, control architectures, and user-facing systems built on normative assumptions about comfort or satisfaction~\cite{becerik2022ten}. While these offer operational clarity, they often abstract away from the socially entangled ways people inhabit built environments. Our studies show that affective experience in smart environments arises through interpretation, bodily attunement, and environmental negotiation. Participants described ``affective drift'' over time, formed ``working theories'' of system behaviour, and expressed tensions between feeling seen and feeling intruded upon. These insights cannot be captured by metrics alone; they require context-sensitive methods attentive to temporality, norms, and somatic cues.

Smart building's knowledge objectives can shift \textit{from seeking generalisable rules to situated understandings}. Rather than focusing on predictability, the field must also embrace epistemic approaches that foreground subjectivity, partiality, and lived experience. This includes qualitative methods like body mapping~\cite{turmo2023towards}, narrative walk-alongs~\cite{resch2015urban}, or temporal mapping~\cite{turmo2025temporal}; interpretive frameworks from affect theory and phenomenology~\cite{lottridge2011affective, fritsch2025sticking}; and hybrid approaches that combine sensing with self-report, conversation, or design probes. The aim is not to dismiss modelling, but to expand what counts as valid knowledge in smart buildings.


% \subsubsection{Evaluation Objectives: Supporting Affective Fit}
% A second prevailing objective has been the pursuit of establishing measurable experiences such as comfort and affect~\cite{bower2019impact, moran2013using}, privacy~\cite{schnadelbach2019adaptive} or agency~\cite{mitchell2020no}. However, our findings suggest that such experiences are not a fixed or stable state. Participants navigated experience through fleeting, embodied impressions of \textit{``affective fit''}---whether a space felt ``good enough'' \textit{at that moment}. These judgments were not absolute, but contingent to bodily changes, social presence, and anticipated demand. 

% This calls for a shift in HBI’s evaluative objectives: \textit{from measuring comfort to cultivating affective fit}. We use \textit{affective fit} to refer to the evolving, somatic, and interpretive sense that a space supports one’s activities, emotions, and intentions, however imperfectly. Unlike comfort, affective fit is not a singular outcome but an ongoing negotiation shaped by bodily state, task demands, social context, and temporal rhythms. Affective fit could be approached as a situated judgment—surfaced through multi-modal reflections (e.g. body maps~\cite{turmo2023towards}, narrative probes~\cite{gaver2004cultural}), relational indicators (e.g. space dwell patterns ~\cite{verma2017studying}), or require the design for new technology that can register stability, drift, and re-alignment over time.

% To support affective fit, HBI must expand and diversify its methodological toolbox. For example, we found body mapping to be especially effective in surfacing tacit spatial judgments among neurodiverse participants in a smart learning space. Yet such approaches are often treated as peripheral or sensitising exercises, rather than core instruments of evaluation. The objective going forward is integrating embodied, interpretive methods with other forms of sensed or behavioural data.  Evaluation, then, must support occupants in articulating and interpreting their own experiences in context, while ensuring those expressions can meaningfully inform building interaction and design.



% \subsubsection{Methodological Objectives: Situated Methodology}
% A longstanding methodological objective in smart building research has been to infer experience through sensed and survey-based data—recording environmental conditions, usage patterns, or subjective ratings to derive generalisable insights. These extractive approaches reflect a broader empirical logic (e.g.~\cite{daissaoui2020iot, dalton2020habitech}), one where knowledge is assumed to arise through the accumulation and analysis of decontextualised data points. Yet such methods often reduce experience to measurable proxies, abstracting away from the temporal, relational, and embodied contexts through which it unfolds.

% In contrast, we draw from interpretive traditions in HCI, particularly those advanced by the affective interaction (AfI) research community~\cite{ahmadpour2025affective}, which foreground the contextual and affective dimensions of lived experience. We find that our use of in-situ walkthroughs and emotion-annotated spatial reflections~\cite{crivellaro2015contesting, mitchell2020no} did not merely capture data—they helped elicit interpretation. Participants expressed discomfort, ambiguity, used spatial metaphors, and drew on “local and vernacular” terms (specific to the building in question) that standardised metrics may miss.

% Unlike acclimatisation-based protocols~\cite{gnecco2025exploring} or approaches that isolate individuals in controlled conditions~\cite{bower2019impact}, our building walk method offered a more dynamic and socially embedded lens on experience. It allowed participants to move, reflect, gesture, and converse in real time—surfacing temporal patterns, transitions, and moments of co-regulation. 

% This calls for a shift in HBI’s methodological objectives: \textit{from sensing and extracting behavioural data to cultivating situated and interpretive depth}. Rather than pursuing ever-finer measurement, the goal is to better attune to how people make sense of their experience in place. This demands methodological pluralism: hybrid approaches that combine sensing with narrative, embodied, and reflexive techniques~\cite{moran2013using, bower2019impact}. Such depth is not about increasing sensor resolution or data volume, but about gaining insight into how people interpret and respond to their environment over time—through bodily cues, social interactions, and evolving spatial judgments.


% \subsubsection{Knowledge Objectives: Developing Situated Understandings}
% A longstanding knowledge objective in smart buildings has been the pursuit of generalisable rules and predictive models---aiming to extract patterns of behaviour or experience that hold across contexts, users, and settings. This approach underpins many optimisation frameworks, control architectures, and even user-facing systems that rely on normative assumptions about comfort, satisfaction, or engagement~\cite{becerik2022ten}. In affective computing and related strands of HCI, this is echoed in efforts to detect and respond to affective states using physiological sensors, facial expressions, or machine learning classifiers~\cite{gao2020n}. While these approaches offer operational clarity, they often reduce experience to abstractable features, limiting attention to the contextually rich, socially entangled, and meaning-laden ways that people actually inhabit built environments.

% Our studies show that affective experience in smart environments emerges through processes of interpretation, bodily attunement, and environmental negotiation. For example, participants articulated ``affective drift'' over time, developed “working theories” of system behaviour, or reported ambiguous tensions between feeling seen and feeling intruded upon. These insights could not be derived from aggregate metrics or discrete affect labels alone—they required reflexive, context-sensitive methods that attended to temporality, social norms, and somatic cues.

% This calls for a shift in HBI’s knowledge objectives: \textit{from seeking generalisable rules to cultivating situated understandings}. Rather than striving for predictability alone, the field must also embrace epistemic approaches that foreground subjectivity, partiality, and lived experience. This includes qualitative methods like body mapping~\cite{turmo2023towards}, narrative walkalongs~\cite{resch2015urban}, or temporal mapping~\cite{turmo2025temporal}; interpretive frameworks from affect theory and phenomenology~\cite{lottridge2011affective, fritsch2025sticking}; and hybrid approaches that combine sensing with self-report, conversation, or design probes. The goal is not to reject formal modelling, but to expand the repertoire of what counts as valid knowledge in HBI. As buildings become more intelligent, the complexity of experience they afford requires equally complex—and situated—ways of knowing.